% !TeX root = ../main.tex
\section{Измерение коэффициента отражения и КСВН фильтра}
Для исследования согласования фильтра нижних частот в полосе пропускания и в полосе заграждения применена измерительная схема на основе направленного ответвителя. Эксперимент базировался на последовательном измерении мощности на выходе ответвленного сигнала для трех характерных состояний нагрузки проходного порта: холостого хода (разомкнутая линия), прецизионной согласованной нагрузки $50\text{ Ом}$ и исследуемого фильтра, нагруженного на сопротивление $50\text{ Ом}$.

Режим разомкнутого порта соответствует полному отражению электромагнитной волны от нагрузки и задает опорный калибровочный уровень отраженной мощности. Измерение с согласованной нагрузкой определяет паразитное просачивание мощности и конечную направленность измерительного ответвителя, ограничивающую динамический диапазон системы на уровне порядка $-31{,}5\text{ дБм}$. Полученные экспериментальные кривые распределения мощности на порте связи приведены на рисунке~\ref{fig:refl_power}.

\begin{figure}[htbp]
    \centering
    \includegraphics[width=0.75\linewidth]{Sparam.png}
    \caption{Спектральное распределение мощности на порте связи (CPL OUT) направленного ответвителя для разомкнутого порта, согласованной нагрузки $50\text{ Ом}$ и фильтра с нагрузкой}
    \label{fig:refl_power}
\end{figure}

Полученные зависимости мощности согласуются с физическими процессами в реактивном фильтре нижних частот. В пределах полосы пропускания на частотах до $400\text{ МГц}$ уровень отраженного от фильтра сигнала совпадает с уровнем согласованной нагрузки ($-31{,}5\text{ дБм}$), что свидетельствует о качественном согласовании входного импеданса фильтра с волновым сопротивлением тракта $50\text{ Ом}$ и беспрепятственном прохождении мощности в нагрузку. В переходной области вблизи частоты среза наблюдается резкий рост мощности, а в полосе заграждения на частотах свыше $650\text{ МГц}$ кривая фильтра сливается с кривой холостого хода (около $-20\text{ дБм}$). Это подтверждает отражательный характер запирания реактивного ФНЧ: электромагнитная волна в полосе заграждения не рассеивается в виде тепла, а отражается обратно к источнику.

Модуль коэффициента отражения фильтра определялся как разность уровней отраженной мощности от фильтра и калибровочного сигнала холостого хода:
\begin{equation}
    |\Gamma(f)|_{\text{дБ}} = P_{\text{filt}}(f) - P_{\text{open}}(f).
\end{equation}

На основе рассчитанного значения модуля коэффициента отражения был вычислен коэффициент стоячей волны по напряжению (КСВН):
\begin{equation}
    \text{КСВН}(f) = \frac{1 + 10^{|\Gamma(f)|_{\text{дБ}}/20}}{1 - 10^{|\Gamma(f)|_{\text{дБ}}/20}}.
\end{equation}

График частотной зависимости КСВН исследуемого фильтра в рабочем диапазоне представлен на рисунке~\ref{fig:ksvn}.

\begin{figure}[htbp]
    \centering
    \includegraphics[width=0.75\linewidth]{КСВН.png}
    \caption{Частотная зависимость коэффициента стоячей волны по напряжению (КСВН) исследуемого фильтра нижних частот}
    \label{fig:ksvn}
\end{figure}

В рабочей полосе пропускания на частотах от $10$ до $430\text{ МГц}$ фильтр демонстрирует устойчивое согласование с измерительным трактом, где значение КСВН не превышает $1{,}6\dots 2{,}0$, что является типичным для реальных устройств с учетом разъемных соединений и паразитных неоднородностей. В переходной области от $450$ до $530\text{ МГц}$ наблюдается интенсивный рост КСВН от $2{,}5$ до $5{,}1$, отражающий стремительное изменение характера входного импеданса фильтра и появление значительной реактивной составляющей.

В полосе заграждения на частотах от $550$ до $600\text{ МГц}$ величина КСВН достигает значений порядка $8{,}5\dots 9{,}3$. При таком уровне КСВН подавляющая часть энергии генератора отражается обратно в тракт, подтверждая высокую эффективность отражательного подавления сигнала фильтром за пределами полосы пропускания.
